\section{Related work}
\label{sec:related_work}

We position this benchmark against six active research directions. In each
case, existing work provides important components, but no single resource
combines causal fault injection, structured state-difference observation,
independent scientific adjudication, calibrated pre-terminal warning, and
frozen transfer splits.

\subsection{Scientific-agent benchmarks}

MLR-Bench evaluates AI agents on open-ended machine learning research tasks
and reports task-completion rates~\citep{mlrbench2025}. SciAgentArena
benchmarks agents across scientific challenges at multiple
scales~\citep{sciagentarena2026}. Both are broad, valuable evaluations of
terminal agent competence. Neither injects controlled faults with known
causal structure, records fine-grained state differences, or measures the
time course of failure. Our benchmark supplements, rather than replaces,
these resources by adding the temporal, causal, and state-difference
dimensions needed to study failure dynamics.

\subsection{Agent-failure taxonomies and diagnosis}

MAST-Data provides more than 1{,}600 annotated multi-agent failure traces
with a descriptive taxonomy~\citep{mast2025}. AgentRx performs critical-step
localization and root-cause diagnosis from execution
trajectories~\citep{agentrx2026}. Aegis combines a failure taxonomy with
targeted intervention suggestions~\citep{aegis2025}. These systems
demonstrate that trajectory-level failure analysis is feasible and
informative. Our contribution differs in three ways: we use governed,
deterministic fault injection rather than post-hoc labeling, so causal
knowledge is established at injection time; we measure failure \emph{dynamics}
(hazard, propagation, recoverability) rather than only presence or class;
and we evaluate calibrated pre-terminal prediction on frozen splits rather
than in-domain diagnosis.

\subsection{Execution monitoring and state difference}

Runtime monitors, state-snapshot systems, and execution tracers are standard
in software engineering and have been applied to agent
systems~\citep{shinn2024}. Our framework standardizes state-difference
observation as a typed contract—recording before/after content hashes, the
semantic type of change, whether the change was expected, and whether
timeout or partial completion was independently reconciled—so that
monitoring data can be used by downstream models without ad-hoc parsing.

\subsection{Calibration, selective prediction, and survival analysis}

Calibrated probability estimation~\citep{guo2017}, selective
prediction~\citep{geifman2017}, and survival or time-to-event
modeling~\citep{cox1972} are individually well studied. Our warning baseline
applies these tools to scientific-agent trajectories, estimating a
time-varying failure hazard, a distribution over causal fault classes, and
an abstention state at each step, with calibration evaluated on dedicated
splits.

\subsection{Fault injection for agents}

Fault injection is established in distributed systems and has been applied
to language-model agents for robustness testing. We extend this tradition
by requiring clean counterfactuals (the same task executed without the
fault), by concealing fault truth from the evaluated policy, by classifying
failed injections separately from successful faults, and by holding out
entire fault families and fault compositions as transfer tests.

\subsection{Recoverability measurement}

Most agent evaluations report final success, but few measure whether an
intervention that arrives after a warning can still restore the trajectory.
Our standardized recovery probe—a fixed, family-appropriate canonical action
applied at early, middle, and late delays—provides a controlled measure of
the closing recoverability window without introducing a learned recovery
policy, which is the subject of the companion method paper.

The novelty claim of this paper is not the existence of a taxonomy, a trace
dataset, a scientific task, or a fault injector. It is the conjunction of
governed causal injection, structured state-difference observation,
independent scientific adjudication, calibrated pre-terminal prediction, and
frozen transfer evaluation—a combination that makes failure dynamics
measurable and the measurement reusable.
